% Generated by the blueprint tool (blueprint-pdf/1) from a blueprint snapshot; edit the
% Markdown sources, not this file.  Compile with `latexmk -pdf` or by running
% `pdflatex` twice (the second run settles the contents and the page references).
% Images, if any, are read from assets/ next to this file.
\documentclass[11pt]{article}
\usepackage[T1]{fontenc}
\usepackage[utf8]{inputenc}
\usepackage{lmodern}
\usepackage{textcomp}
\usepackage[a4paper,margin=2.5cm,headheight=14pt]{geometry}
\usepackage{amsmath,amssymb,amsthm}
\usepackage{amscd}  % KaTeX's CD diagrams; part of amsmath, so always there
\usepackage{mathtools}
\usepackage{bm}
\usepackage{graphicx}
\usepackage{xcolor}
\usepackage{array,booktabs,tabularx}
\usepackage{fancyvrb}
\usepackage[normalem]{ulem}
\usepackage{microtype}
% Optional, and only with their fonts: a small TeX installation has the
% packages without the fonts, which fails at the first formula.
\IfFileExists{mathrsfs.sty}{\IfFileExists{rsfs10.tfm}{\usepackage{mathrsfs}}{}}{}
\IfFileExists{stmaryrd.sty}{\IfFileExists{stmary10.tfm}{\usepackage{stmaryrd}}{}}{}
\IfFileExists{cancel.sty}{\usepackage{cancel}}{}
\usepackage[hyphens]{url}
\usepackage[pdfusetitle,bookmarksnumbered=false]{hyperref}
\usepackage{bookmark}
\definecolor{bplink}{HTML}{1F4E9C}
\definecolor{bpmuted}{HTML}{5F676E}
\hypersetup{colorlinks=true,linkcolor=bplink,urlcolor=bplink,citecolor=bplink}
\setlength{\parindent}{0pt}
\setlength{\parskip}{0.55em plus 0.2em minus 0.1em}
% Lists as tight as the page's.
\makeatletter
\def\@listi{\leftmargin\leftmargini\parsep 0pt\topsep 0.3em plus 0.1em\itemsep 0.1em plus 0.1em}
\let\@listI\@listi
\def\@listii{\leftmargin\leftmarginii\labelwidth\leftmarginii\advance\labelwidth-\labelsep\parsep 0pt\topsep 0.1em\itemsep 0.1em}
% An ordered list that does not start at 1.
\newcommand{\bpstartat}[1]{\setcounter{\@enumctr}{\numexpr#1-1\relax}}
\makeatother
\setcounter{tocdepth}{3}
\emergencystretch=3em
\allowdisplaybreaks

% A numbered statement, set as a paper sets one: "Theorem 2.1 (Title)."
% #1 is the lead ("Theorem~2.1"), #2 what follows it in roman.
\newenvironment{bpstatement}[2]{\par\addvspace{0.6\baselineskip}\noindent\textbf{#1}#2\textbf{.}\enspace\ignorespaces}
  {\par\addvspace{0.4\baselineskip}}
\newenvironment{bpstatementit}[2]{\begin{bpstatement}{#1}{#2}\itshape\ignorespaces}{\end{bpstatement}}
% A step: a relation with prose of its own, read under the object it starts from.
\newenvironment{bpstep}[1]{\par\addvspace{0.3\baselineskip}\begin{list}{}{\leftmargin=1.6em\rightmargin=0pt\topsep=0pt}\item\relax
  {\small\color{bpmuted}#1}\par\nopagebreak}{\end{list}\par}
% Notes about where an object is written out.
\newenvironment{bpnote}{\par\small\color{bpmuted}\itshape}{\par}
\newcommand{\bpomitted}[1]{\par{\small\color{bpmuted}\itshape #1\par}}
\newcommand{\bpbodyheading}[1]{\par\addvspace{0.4\baselineskip}\noindent\textbf{#1}\par\nopagebreak}
\newcommand{\bprule}{\par\medskip\noindent\hfil\rule{0.3\linewidth}{0.4pt}\hfil\par\medskip}
\newcommand{\bpboundary}[1]{\par\nopagebreak{\small\color{bpmuted}#1\par}\nopagebreak}
% A reference to something this export leaves out, and one that resolves nowhere.
\newcommand{\bpout}[1]{#1\textsuperscript{\hyperlink{bp-outside}{\ensuremath{\ast}}}}
\newcommand{\bpbroken}[1]{\textcolor{red!70!black}{[#1?]}}
\newcommand{\bpslug}[1]{\texttt{#1}}
% A bibliography entry: label, anchor, entry; the label hangs in the margin.
\newcommand{\bpbibitem}[3]{\par\setbox0\hbox{[#1]\enspace}\ifdim\wd0<5em\setbox0\hbox to 5em{[#1]\hfil}\fi
  \hangindent=\wd0 \hangafter=1 \noindent\box0\phantomsection\label{#2}#3\par\smallskip}
% A character the export knows no way to set.
\newcommand{\bpglyph}[1]{\fbox{\scriptsize #1}}
% Images, at their own size unless they are wider than the line.
\newsavebox{\bpimagebox}
\newcommand{\bpimage}[1]{\sbox{\bpimagebox}{\includegraphics{#1}}%
  \ifdim\wd\bpimagebox>\linewidth\resizebox{\linewidth}{!}{\usebox{\bpimagebox}}\else\usebox{\bpimagebox}\fi}
\newcommand{\bpmissingimage}[1]{\fbox{\parbox{0.8\linewidth}{\centering\small\color{bpmuted}Image not included: #1}}}
\DefineVerbatimEnvironment{bpverbatim}{Verbatim}{fontsize=\small,xleftmargin=1em}
\renewcommand{\qedsymbol}{\ensuremath{\blacksquare}}
\hypersetup{pdftitle={Weil Conjectures — 13 Weights},pdfsubject={Chapter 13 Weights; all levels},pdfcreator={blueprint}}
% KaTeX commands the website knows
\providecommand{\R}{\mathbb{R}}\providecommand{\N}{\mathbb{N}}\providecommand{\Z}{\mathbb{Z}}
\providecommand{\Q}{\mathbb{Q}}\providecommand{\C}{\mathbb{C}}
\providecommand{\Reals}{\mathbb{R}}\providecommand{\reals}{\mathbb{R}}\providecommand{\Complex}{\mathbb{C}}
\providecommand{\cnums}{\mathbb{C}}\providecommand{\natnums}{\mathbb{N}}\providecommand{\Bbb}[1]{\mathbb{#1}}
\providecommand{\bold}[1]{\mathbf{#1}}\providecommand{\frak}[1]{\mathfrak{#1}}\providecommand{\mathbbm}[1]{\mathbb{#1}}
\providecommand{\lang}{\langle}\providecommand{\rang}{\rangle}\providecommand{\bra}[1]{\langle #1\rvert}
\providecommand{\ket}[1]{\lvert #1\rangle}\providecommand{\braket}[1]{\langle #1\rangle}
\providecommand{\set}[1]{\{#1\}}\providecommand{\Set}[1]{\left\{#1\right\}}
\providecommand{\infin}{\infty}\providecommand{\isin}{\in}\providecommand{\plusmn}{\pm}\providecommand{\sub}{\subset}
\providecommand{\sube}{\subseteq}\providecommand{\supe}{\supseteq}\providecommand{\exist}{\exists}
\providecommand{\bull}{\bullet}\providecommand{\larr}{\leftarrow}\providecommand{\rarr}{\rightarrow}
\providecommand{\lrarr}{\leftrightarrow}\providecommand{\Larr}{\Leftarrow}\providecommand{\Rarr}{\Rightarrow}
\providecommand{\Lrarr}{\Leftrightarrow}\providecommand{\uarr}{\uparrow}\providecommand{\darr}{\downarrow}
\providecommand{\harr}{\leftrightarrow}\providecommand{\lArr}{\Leftarrow}\providecommand{\rArr}{\Rightarrow}
\providecommand{\hArr}{\Leftrightarrow}\providecommand{\alef}{\aleph}\providecommand{\alefsym}{\aleph}
\providecommand{\real}{\Re}\providecommand{\image}{\Im}\providecommand{\weierp}{\wp}
\providecommand{\Alpha}{\mathrm{A}}\providecommand{\Beta}{\mathrm{B}}\providecommand{\Epsilon}{\mathrm{E}}
\providecommand{\Zeta}{\mathrm{Z}}\providecommand{\Eta}{\mathrm{H}}\providecommand{\Iota}{\mathrm{I}}
\providecommand{\Kappa}{\mathrm{K}}\providecommand{\Mu}{\mathrm{M}}\providecommand{\Nu}{\mathrm{N}}
\providecommand{\Omicron}{\mathrm{O}}\providecommand{\Rho}{\mathrm{P}}\providecommand{\Tau}{\mathrm{T}}
\providecommand{\Chi}{\mathrm{X}}\providecommand{\omicron}{o}\providecommand{\gt}{>}\providecommand{\lt}{<}
\providecommand{\And}{\&}\providecommand{\hdashline}{\hline}\providecommand{\operatornamewithlimits}[1]{\operatorname*{#1}}
\providecommand{\llbracket}{[\![}\providecommand{\rrbracket}{]\!]}\providecommand{\cancel}[1]{#1}
\providecommand{\mathscr}{\mathcal}
% the characters beyond ASCII this document uses
\DeclareUnicodeCharacter{2160}{\textup{I}}
\DeclareUnicodeCharacter{2161}{\textup{II}}
\DeclareUnicodeCharacter{2162}{\textup{III}}
% the project's macros ([katex.macros] in blueprint.toml)
\providecommand{\A}{}\renewcommand{\A}{\mathbf A}
\providecommand{\AEn}{}\renewcommand{\AEn}{A_{E,n}}
\providecommand{\Aaf}{}\renewcommand{\Aaf}{\mathbf A}
\providecommand{\Ab}{}\renewcommand{\Ab}{\mathsf{Ab}}
\providecommand{\AffSch}{}\renewcommand{\AffSch}{\operatorname{AffSch}}
\providecommand{\Aone}{}\renewcommand{\Aone}{\mathbf A^1}
\providecommand{\Bl}{}\renewcommand{\Bl}{\operatorname{Bl}}
\providecommand{\Br}{}\renewcommand{\Br}{\mathrm{Br}}
\providecommand{\C}{}\renewcommand{\C}{\mathbb C}
\providecommand{\Cech}{}\renewcommand{\Cech}{\check{H}}
\providecommand{\Coext}{}\renewcommand{\Coext}{\operatorname{Coext}}
\providecommand{\Cone}{}\renewcommand{\Cone}{\operatorname{Cone}}
\providecommand{\Cons}{}\renewcommand{\Cons}{\operatorname{Cons}}
\providecommand{\Cref}{}\renewcommand{\Cref}[1]{\text{#1}}
\providecommand{\D}{}\renewcommand{\D}{D}
\providecommand{\DLAFT}{}\renewcommand{\DLAFT}{\operatorname{FT}}
\providecommand{\DLAFr}{}\renewcommand{\DLAFr}{\operatorname{Fr}}
\providecommand{\DLAQ}{}\renewcommand{\DLAQ}{{\overline{\mathbf Q}_\ell}}
\providecommand{\DLRQ}{}\renewcommand{\DLRQ}{{\overline{\mathbf Q}_\ell}}
\providecommand{\Dbc}{}\renewcommand{\Dbc}{D^b_c}
\providecommand{\Dc}{}\renewcommand{\Dc}{D_{\mathrm{cons}}}
\providecommand{\Dctf}{}\renewcommand{\Dctf}{D^b_{ctf}}
\providecommand{\Dic}{}\renewcommand{\Dic}{D_{\IndCons}}
\providecommand{\Dplus}{}\renewcommand{\Dplus}{D^{+}}
\providecommand{\End}{}\renewcommand{\End}{\operatorname{End}}
\providecommand{\Et}{}\renewcommand{\Et}{\mathrm{\acute{e}t}}
\providecommand{\Ex}{}\renewcommand{\Ex}{\mathrm{Ex}}
\providecommand{\Ext}{}\renewcommand{\Ext}{\operatorname{Ext}}
\providecommand{\F}{}\renewcommand{\F}{\mathbb F}
\providecommand{\FT}{}\renewcommand{\FT}{\operatorname{FT}}
\providecommand{\Fbar}{}\renewcommand{\Fbar}{{\overline{\mathbf F}_q}}
\providecommand{\Fell}{}\renewcommand{\Fell}{\mathbf F_\ell}
\providecommand{\Fib}{}\renewcommand{\Fib}{\operatorname{Fib}}
\providecommand{\Fix}{}\renewcommand{\Fix}{\operatorname{Fix}}
\providecommand{\Fp}{}\renewcommand{\Fp}{{\mathbf F_{p}}}
\providecommand{\Fq}{}\renewcommand{\Fq}{{\mathbf F_q}}
\providecommand{\Fqn}{}\renewcommand{\Fqn}{{\mathbf F_{q^{n}}}}
\providecommand{\Frob}{}\renewcommand{\Frob}{\operatorname{Frob}}
\providecommand{\GL}{}\renewcommand{\GL}{\operatorname{GL}}
\providecommand{\Gal}{}\renewcommand{\Gal}{\operatorname{Gal}}
\providecommand{\Hens}{}\renewcommand{\Hens}{\operatorname{Hens}}
\providecommand{\Hilb}{}\renewcommand{\Hilb}{\operatorname{Hilb}}
\providecommand{\Hom}{}\renewcommand{\Hom}{\operatorname{Hom}}
\providecommand{\IndCons}{}\renewcommand{\IndCons}{\operatorname{IndCons}}
\providecommand{\Kar}{}\renewcommand{\Kar}{\operatorname{Kar}}
\providecommand{\LT}{}\renewcommand{\LT}{\operatorname{LT}}
\providecommand{\Lie}{}\renewcommand{\Lie}{\operatorname{Lie}}
\providecommand{\Loc}{}\renewcommand{\Loc}{\operatorname{Loc}}
\providecommand{\Mod}{}\renewcommand{\Mod}{\mathsf{Mod}}
\providecommand{\N}{}\renewcommand{\N}{\mathbb{N}}
\providecommand{\Needspace}{}\renewcommand{\Needspace}{\par}
\providecommand{\O}{}\renewcommand{\O}{{\mathcal{O}}}
\providecommand{\Ob}{}\renewcommand{\Ob}{\mathrm{Ob}}
\providecommand{\Oe}{}\renewcommand{\Oe}{\mathcal O_E}
\providecommand{\PP}{}\renewcommand{\PP}{\mathbf{P}}
\providecommand{\Perf}{}\renewcommand{\Perf}{\operatorname{Perf}}
\providecommand{\Pic}{}\renewcommand{\Pic}{\mathrm{Pic}}
\providecommand{\Pm}{}\renewcommand{\Pm}{\mathbf P^{m}}
\providecommand{\Pone}{}\renewcommand{\Pone}{\mathbf P^{1}}
\providecommand{\Pqn}{}\renewcommand{\Pqn}{\mathbf P^{n}}
\providecommand{\ProFin}{}\renewcommand{\ProFin}{\operatorname{ProFin}}
\providecommand{\Proet}{}\renewcommand{\Proet}{\mathrm{pro\acute{e}t}}
\providecommand{\Q}{}\renewcommand{\Q}{\mathbb Q}
\providecommand{\Qb}{}\renewcommand{\Qb}{\overline{\mathbf Q}_{\ell}}
\providecommand{\Qbar}{}\renewcommand{\Qbar}{\overline{\mathbf Q}}
\providecommand{\QbarX}{}\renewcommand{\QbarX}{\overline{\mathbf Q}_{\ell,X}}
\providecommand{\Qbarl}{}\renewcommand{\Qbarl}{{\overline{\mathbf Q}_{\ell}}}
\providecommand{\Qell}{}\renewcommand{\Qell}{{\mathbf{Q}_{\ell}}}
\providecommand{\Ql}{}\renewcommand{\Ql}{\mathbf Q_{\ell}}
\providecommand{\Qlb}{}\renewcommand{\Qlb}{{\overline{\mathbf Q}_\ell}}
\providecommand{\Qlbar}{}\renewcommand{\Qlbar}{\overline{\mathbf Q}_\ell}
\providecommand{\R}{}\renewcommand{\R}{\mathbb{R}}
\providecommand{\RG}{}\renewcommand{\RG}{\operatorname{\mathrm{R}\Gamma}}
\providecommand{\RGamma}{}\renewcommand{\RGamma}{R\Gamma}
\providecommand{\RGammac}{}\renewcommand{\RGammac}{R\Gamma_c}
\providecommand{\RHom}{}\renewcommand{\RHom}{R\mathcal{H}om}
\providecommand{\RHomX}{}\renewcommand{\RHomX}{R\mathcal H om_X}
\providecommand{\Rat}{}\renewcommand{\Rat}{\operatorname{IsRat}}
\providecommand{\Res}{}\renewcommand{\Res}{\operatorname{Res}}
\providecommand{\Rlim}{}\renewcommand{\Rlim}{R\!\varprojlim}
\providecommand{\SIm}{}\renewcommand{\SIm}{\operatorname{SIm}}
\providecommand{\Sch}{}\renewcommand{\Sch}{\mathsf{Sch}}
\providecommand{\Sh}{}\renewcommand{\Sh}{\operatorname{Sh}}
\providecommand{\Span}{}\renewcommand{\Span}{\operatorname{span}}
\providecommand{\Spec}{}\renewcommand{\Spec}{\operatorname{Spec}}
\providecommand{\Tr}{}\renewcommand{\Tr}{\operatorname{Tr}}
\providecommand{\Trsm}{}\renewcommand{\Trsm}{\operatorname{Tr}^{\mathrm{sm}}}
\providecommand{\Wt}{}\renewcommand{\Wt}{\operatorname{Wt}}
\providecommand{\Z}{}\renewcommand{\Z}{\mathbb{Z}}
\providecommand{\Zar}{}\renewcommand{\Zar}{\mathrm{zar}}
\providecommand{\Zell}{}\renewcommand{\Zell}{{\mathbf{Z}_{\ell}}}
\providecommand{\Zl}{}\renewcommand{\Zl}{{\mathbf Z_\ell}}
\providecommand{\affet}{}\renewcommand{\affet}[1]{{#1}^{\mathrm{aff}}_{\mathrm{ét}}}
\providecommand{\affproet}{}\renewcommand{\affproet}[1]{{#1}^{\mathrm{aff}}_{\mathrm{proét}}}
\providecommand{\boxt}{}\renewcommand{\boxt}{\boxtimes_A^{\mathrm L}}
\providecommand{\cF}{}\renewcommand{\cF}{\mathcal{F}}
\providecommand{\cHom}{}\renewcommand{\cHom}{R\mathcal{H}\!om}
\providecommand{\cI}{}\renewcommand{\cI}{\mathcal{I}}
\providecommand{\cL}{}\renewcommand{\cL}{\mathcal L}
\providecommand{\cM}{}\renewcommand{\cM}{\mathcal M}
\providecommand{\cN}{}\renewcommand{\cN}{\mathcal N}
\providecommand{\cO}{}\renewcommand{\cO}{\mathcal O}
\providecommand{\cT}{}\renewcommand{\cT}{\mathcal T}
\providecommand{\checked}{}\renewcommand{\checked}{\color{red}\aftergroup\restorecolor}
\providecommand{\coeff}{}\renewcommand{\coeff}{\operatorname{coeff}}
\providecommand{\coker}{}\renewcommand{\coker}{\operatorname{coker}}
\providecommand{\colim}{}\renewcommand{\colim}{\operatorname{colim}}
\providecommand{\constructedby}{}\renewcommand{\constructedby}[1]{\par\smallskip\noindent\textcolor{purple!60!black}{\textsf{Constructed by:} }
  \footnotesize\sloppy\nolinkurl{#1}\par}
\providecommand{\deepresult}{}\renewcommand{\deepresult}[1]{\par\smallskip\noindent\textcolor{orange!45!black}{
  \textsf{Blueprint role: named deep theorem node; see the cluster proof architecture.}}
  \par\noindent\textcolor{orange!45!black}{\textsf{Mathematical source:} }#1\par}
\providecommand{\dt}{}\renewcommand{\dt}{\mathbin{\otimes^{\mathbf L}}}
\providecommand{\dw}{}\renewcommand{\dw}{\operatorname{dw}}
\providecommand{\et}{}\renewcommand{\et}[1]{{#1}_{\mathrm{ét}}}
\providecommand{\ev}{}\renewcommand{\ev}{\operatorname{ev}}
\providecommand{\fib}{}\renewcommand{\fib}{\operatorname{fib}}
\providecommand{\formalizationnote}{}\renewcommand{\formalizationnote}[1]{\begin{quote}\small\raggedright\textbf{Formalization note.} #1\end{quote}}
\providecommand{\fpqc}{}\renewcommand{\fpqc}{X_{\mathrm{fpqc}}}
\providecommand{\gammacont}{}\renewcommand{\gammacont}{\Gamma_{\mathrm{cont}}}
\providecommand{\gm}{}\renewcommand{\gm}{\mathbf G_m}
\providecommand{\hocolim}{}\renewcommand{\hocolim}{\operatorname{hocolim}}
\providecommand{\holim}{}\renewcommand{\holim}{\operatorname{holim}}
\providecommand{\hrB}{}\renewcommand{\hrB}{\@ifnextchar]{\hrEnd}{\hrStep}}
\providecommand{\hrEnd}{}\renewcommand{\hrEnd}{\@firstoftwo{}}
\providecommand{\hrStep}{}\renewcommand{\hrStep}{\@firstoftwo{\hrB}}
\providecommand{\id}{}\renewcommand{\id}{\operatorname{id}}
\providecommand{\kunFl}{}\renewcommand{\kunFl}{\mathbf F_\ell}
\providecommand{\kunZn}{}\renewcommand{\kunZn}{\mathbf Z/n\mathbf Z}
\providecommand{\kuncd}{}\renewcommand{\kuncd}{\operatorname{cd}}
\providecommand{\kundim}{}\renewcommand{\kundim}{\dim}
\providecommand{\kunotimes}{}\renewcommand{\kunotimes}[1]{\otimes^{\mathbf L}_{#1}}
\providecommand{\kunsupp}{}\renewcommand{\kunsupp}{\operatorname{Supp}}
\providecommand{\lotimes}{}\renewcommand{\lotimes}{\otimes^{\mathrm L}}
\providecommand{\m}{}\renewcommand{\m}{{\mathfrak{m}}}
\providecommand{\n}{}\renewcommand{\n}{{\mathfrak{n}}}
\providecommand{\nagataGap}{}\renewcommand{\nagataGap}[1]{\par\noindent\textbf{Formalization gap.} #1\par}
\providecommand{\openobligation}{}\renewcommand{\openobligation}[1]{\par\smallskip\noindent\textcolor{red!65!black}{\textsf{Open obligation:} }
  \footnotesize\sloppy\nolinkurl{#1}\par}
\providecommand{\p}{}\renewcommand{\p}{{\mathfrak{p}}}
\providecommand{\pr}{}\renewcommand{\pr}{\operatorname{pr}}
\providecommand{\preshv}{}\renewcommand{\preshv}[1]{\mathrm{PreShv}(#1)}
\providecommand{\pro}{}\renewcommand{\pro}[1]{\mathrm{Pro}(#1)}
\providecommand{\proet}{}\renewcommand{\proet}{X_{\mathrm{proét}}}
\providecommand{\q}{}\renewcommand{\q}{{\mathfrak{q}}}
\providecommand{\rad}{}\renewcommand{\rad}{\operatorname{rad}}
\providecommand{\rank}{}\renewcommand{\rank}{\operatorname{rank}}
\providecommand{\red}{}\renewcommand{\red}{\color{red}}
\providecommand{\restorecolor}{}\renewcommand{\restorecolor}{\color{black}}
\providecommand{\retiredobligation}{}\renewcommand{\retiredobligation}[1]{\par\smallskip\noindent\textcolor{red!75!black}{\textsf{Retired invalid route:} }
  \footnotesize\sloppy\nolinkurl{#1}\par}
\providecommand{\rk}{}\renewcommand{\rk}{\operatorname{rk}}
\providecommand{\sep}{}\renewcommand{\sep}{\textnormal{sep}}
\providecommand{\shv}{}\renewcommand{\shv}[1]{\mathrm{Shv}(#1)}
\providecommand{\spec}{}\renewcommand{\spec}[1]{\mathrm{Spec}(#1)}
\providecommand{\stacks}{}\renewcommand{\stacks}[1]{\href{https://stacks.math.columbia.edu/tag/#1}{\textnormal{Stacks, Tag~#1}}}
\providecommand{\stacksproject}{}\renewcommand{\stacksproject}[1]{{\cite[\href{https://stacks.math.columbia.edu/tag/#1}{Tag #1}]{stacks-project}}}
\providecommand{\statusnote}{}\renewcommand{\statusnote}[1]{\begin{quote}\small\raggedright\textbf{Status.} #1\end{quote}}
\providecommand{\tildeideal}{}\renewcommand{\tildeideal}[2]{\widetilde{(#1, #2)}}
\providecommand{\tildering}{}\renewcommand{\tildering}[3]{{#1}_{\widetilde{(#2, #3)}}}
\providecommand{\tr}{}\renewcommand{\tr}{\operatorname{tr}}
\providecommand{\wt}{}\renewcommand{\wt}{\operatorname{wt}}

\begin{document}
\pagestyle{plain}
\begin{center}
{\LARGE\bfseries Weil Conjectures\par}
\vspace{0.7em}
{\large Chapter 13: Weights\par}
\vspace{0.5em}
{\small\color{bpmuted}Snapshot 2026-09-26, revision 8e4cb88e2a\par}
\end{center}
\vspace{0.5em}
{\small\tableofcontents}
\bigskip\hrule\bigskip

% object: chap-Weights (blueprint/chap-Weights/_section.md)
\section*{\texorpdfstring{13\quad Weights}{13 Weights}}
\phantomsection\label{bp1}%
\addcontentsline{toc}{section}{\texorpdfstring{13\quad Weights}{13 Weights}}
% text: prose-13-01
\phantomsection\label{bp2}%
For the elliptic curve \(y^2=x^3-x\) over \(\mathbf F_3\) of the introduction, Frobenius acts on \(H^1\) with eigenvalues \(\pm i\sqrt3\). The Riemann hypothesis for this curve is then the statement \(|\pm i\sqrt3|=3^{1/2}\). But the eigenvalues live in \(\overline{\mathbf Q}_\ell\), and an \(\ell\)-adic number has no complex absolute value until we choose a field embedding \(\iota:\overline{\mathbf Q}_\ell\hookrightarrow\mathbf C\). Since no choice is preferred, we fix one \(\iota\) throughout, and the final theorems hold for every choice.

An eigenvalue \(\alpha\) of Frobenius at a closed point of degree \(n\) has weight \(w\) if \(|\iota(\alpha)|=q^{nw/2}\). The factor \(n\) makes the weight independent of the field over which we look at the point: over \(\mathbf F_{q^{mn}}\) the eigenvalue becomes \(\alpha^m\) and \(q^n\) becomes \(q^{mn}\). For example, on \(\mathbf P^1\) Frobenius acts on \(H^0\) by \(1\), of weight \(0\), and on \(H^2=\overline{\mathbf Q}_\ell(-1)\) by \(q\), of weight \(2\). Similarly, the eigenvalues on \(H^1\) of our elliptic curve have weight \(1\). Weil's purity axiom says that this pattern persists: \(H^i\) of a smooth proper variety is pure of weight \(i\). The name comes from Hodge theory, where \(H^i\) of a compact Kähler manifold carries a Hodge structure of weight \(i\).

The proof cannot stay with smooth proper varieties. It fibres them in curves, and the cohomology of the fibres forms sheaves on the base. So we need weights for sheaves, and for cohomology that is not pure. Without smoothness and properness only an upper bound survives. For example, let \(\mathbf G_m=\mathbf P^1\setminus\{0,\infty\}\). The localization sequence gives \(H^1_c(\mathbf G_m,\overline{\mathbf Q}_\ell)=\overline{\mathbf Q}_\ell\) with trivial Frobenius and \(H^2_c(\mathbf G_m,\overline{\mathbf Q}_\ell)=\overline{\mathbf Q}_\ell(-1)\). This agrees with the trace formula \(q-1=\#\mathbf G_m(\mathbf F_q)\). Thus \(H^1_c\) has weight \(0\), strictly less than \(1\).

A sheaf is pure of weight \(w\) if all its eigenvalues at all closed points have weight \(w\). It is mixed if it has a finite filtration whose quotients are pure. A bounded complex \(K\) has weights at most \(w\) if each \(\mathcal H^i(K)\) is mixed with all weights at most \(w+i\) (Definition \hyperref[bp7]{13.1}). Here the shift by \(i\) puts the constant sheaf and its cohomology into one statement: \(R\Gamma_c(X,\overline{\mathbf Q}_\ell)\) has weights at most \(0\) exactly when \(H^i_c\) has weights at most \(i\), as on \(\mathbf G_m\). In these terms the Deligne–Laumon theorem (Theorem \hyperref[bp36]{13.11}) says that \(Rf_!\) does not increase weights.

Two warnings will matter later. The first is that weights depend on \(\iota\). For example, consider the rank-one Weil sheaf on \(\operatorname{Spec}\mathbf F_q\) on which Frobenius acts by \(1+\sqrt2\in\overline{\mathbf Q}_\ell\). It has weight \(2\log_q(1+\sqrt2)>0\) for an embedding with \(\iota(\sqrt2)=\sqrt2\), and weight \(2\log_q(\sqrt2-1)=-2\log_q(1+\sqrt2)\) for the other one. For this reason Theorem \bpout{19.7} needs the Deligne–Laumon theorem for every \(\iota\): only then do the reciprocal roots of \(P_i\) have absolute value \(q^{i/2}\) under every complex embedding.

The second warning is that being mixed is not a pointwise condition: it asks for a filtration with pure quotients, and such a filtration has to be constructed. Most of Chapter \bpout{18} is spent constructing one for direct images.

We define weights on Weil sheaves, sheaves on \(X\) with a Frobenius structure, rather than on sheaves on \(X_0\). The reason is that normalizing weights requires twisting by the geometrically constant rank-one Weil sheaves \(\mathcal B_b\), on which a degree-one Frobenius acts by a scalar \(b\in\overline{\mathbf Q}_\ell^\times\). These do not come from \(X_0\) when \(b\) is not an \(\ell\)-adic unit.

% object: sec-Weil-sheaves-and-weights (blueprint/chap-Weights/sec-Weil-sheaves-and-weights/_section.md), content hidden
\subsection*{\texorpdfstring{13.1\quad Weil sheaves and weights}{13.1 Weil sheaves and weights}}
\phantomsection\label{bp3}%
\addcontentsline{toc}{subsection}{\texorpdfstring{13.1\quad Weil sheaves and weights}{13.1 Weil sheaves and weights}}
% text: prose-13-02
\phantomsection\label{bp4}%
Weights are defined by the absolute values of the Frobenius eigenvalues on stalks at closed points, so we first fix how Frobenius acts on sheaves.

On a smooth curve \(U_0\) we want to treat lisse Weil sheaves as representations, just as a lisse sheaf on \(U_0\) is a continuous representation of \(\pi_1(U_0)\). However, the profinite group \(\pi_1(U_0)\) is an extension of \(\widehat{\mathbf Z}\) by \(\pi_1(U)\), and it has too few representations for this. Indeed, a continuous representation of a profinite group preserves a lattice, so a lift \(s\) of Frobenius acts with eigenvalues that are \(\ell\)-adic units. Hence the rank-one Weil sheaf \(\mathcal B_\ell\), on which \(s\) acts by \(\ell\) and \(\pi_1(U)\) acts trivially, is not a representation of \(\pi_1(U_0)\). Concretely, \(n\mapsto\ell^n\) does not extend continuously to \(\widehat{\mathbf Z}\): the integers \(m!\) tend to \(0\) in \(\widehat{\mathbf Z}\), whereas \(\ell^{m!}\) tends to \(0\) in \(\mathbf Z_\ell\) and not to \(1\).

The Weil group keeps \(\pi_1(U)\) with its topology and replaces \(\widehat{\mathbf Z}\) by the discrete group \(\mathbf Z\) of integer powers of Frobenius. Continuity then constrains only the geometric part. These representations of the Weil group are the language of Theorem \bpout{14.5} and of the monodromy groups of Section \bpout{14.2}.

% text: situation-weil
\phantomsection\label{bp5}%
\textbf{Convention.} Put \(k_0=\mathbf F_q\), \(k=\overline{\mathbf F}_q\), and fix a prime \(\ell\nmid q\) and a field embedding \(\iota:\overline{\mathbf Q}_\ell\hookrightarrow\mathbf C\). For \(X_0/k_0\) separated of finite type, we write \(X=X_0\times_{k_0}k\), and for a closed point \(x\) of \(X_0\) write \(\deg(x)=[\kappa(x):k_0]\) for its degree and \(q_x=q^{\deg(x)}=\#\kappa(x)\). Moreover, geometric Frobenius, Weil complexes and their Frobenius operators are those of Section \bpout{10.5}. As there, a Weil sheaf is a Weil complex concentrated in degree zero, and the six operations carry Weil structures to Weil structures. In this Part Weil sheaves and complexes have coefficients in \(\overline{\mathbf Q}_\ell\); a Weil complex \(K\) with coefficients in a finite extension \(E\subset\overline{\mathbf Q}_\ell\) of \(\mathbf Q_\ell\) is called pure, mixed, or of weights at most \(w\) if \(K\otimes_E\overline{\mathbf Q}_\ell\) is (see Lemma \hyperref[bp9]{13.2} and the remark after it for why scalars must be extended).

On a connected normal scheme of finite type over a field, we identify a lisse \(E\)-sheaf with the corresponding continuous representation of the profinite fundamental group, by Proposition \hyperref[bp32]{13.10}, proved in Section \hyperref[bp24]{13.2} below. This covers all arguments with representations below, since they take place on smooth schemes. Moreover, a lisse Weil \(\overline{\mathbf Q}_\ell\)-sheaf has a lisse Weil model over some finite extension \(E/\mathbf Q_\ell\), and two models become identified after a finite extension of scalars. This model includes the Weil isomorphism, which need not preserve a geometric \(\mathcal O_E\)-lattice.

For a geometrically connected smooth curve \(U_0\), put \(U=U_0\times_{k_0}k\) and choose a geometric point \(\bar\eta\) of \(U\). We then write \(\pi_1(U)\) and \(\pi_1(U_0)\) for \(\pi_1(U,\bar\eta)\) and \(\pi_1(U_0,\bar\eta)\).  The inverse image of the integer powers of geometric Frobenius in \(\pi_1(U_0)\) is its Weil group \(W(U_0)\).  It contains the geometric fundamental group, with an exact sequence

\[
1\longrightarrow \pi_1(U)\longrightarrow W(U_0)
\xrightarrow{\deg}\mathbf Z\longrightarrow1.
\]

A degree-one lift \(s\) splits this sequence as a group extension. In these terms, lisse Weil sheaves are the same as representations \(\rho\) of \(W(U_0)\), defined over a finite extension of \(\mathbf Q_\ell\), whose restriction to \(\pi_1(U)\) is continuous. For \(x\in|U_0|\), the local Frobenius \(\operatorname{Frob}_x\) at \(x\) (Theorem \bpout{12.30}) is a conjugacy class of degree \(\deg(x)\).

% text: prose-13-03
\phantomsection\label{bp6}%
Every sheaf defined over \(k_0\) is a Weil sheaf. By means of the Weil isomorphism, Frobenius acts on the stalk at a geometric point over each closed point \(x\), so we can speak of the eigenvalues of \(\operatorname{Frob}_x\). In the description by representations the Weil isomorphism is transport by \(s\), and conversely it determines the action of \(s\).

The definition below contains two kinds of bounds. The first, being punctually of weights \(\le w\), bounds the eigenvalues at each closed point; it is what the curve case yields on each stalk of a direct image. But it cannot serve as input for the curve case, which needs pure lisse sheaves. The second, mixedness, has the opposite virtue: a mixed sheaf can be treated one pure quotient at a time, and for a mixed sheaf the bound on weights is again a pointwise condition (Lemma \hyperref[bp12]{13.3}). In Chapter \bpout{18} we use this to turn the stalkwise bounds into the bound of Theorem \hyperref[bp36]{13.11}.

% object: definition-weights (blueprint/chap-Weights/sec-Weil-sheaves-and-weights/definition-weights/_section.md)
\phantomsection\label{bp7}%
\begin{bpstatement}{Definition~13.1}{ (Weights, purity and mixedness)}
Let \(K\) be a constructible Weil sheaf on \(X\), let \(x\) be a closed point of \(X_0\), and let \(\alpha\) be an eigenvalue of \(\operatorname{Frob}_x\) on \(K_{\bar x}\). We say \(\alpha\) has \emph{\(\iota\)-weight \(w\)} if \(|\iota(\alpha)|=q_x^{w/2}\).

\begin{enumerate}
\item \(K\) is \emph{punctually \(\iota\)-pure of weight \(w\)}, or simply \emph{pure of weight \(w\)}, if every such eigenvalue, at every closed point, has \(\iota\)-weight \(w\).
\item \(K\) is \emph{punctually of weights \(\le w\)} if \(|\iota(\alpha)|\le q_x^{w/2}\) for every such eigenvalue at every closed point.
\item \(K\) is \emph{\(\iota\)-mixed} if it has a finite filtration by Weil subsheaves whose nonzero quotients are pure, of possibly different real weights. It is \emph{mixed of weights at most \(w\)}, or has \emph{weights at most \(w\)}, if it has such a filtration whose nonzero quotients have weights at most \(w\).
\end{enumerate}

A bounded constructible Weil complex \(K\) has \emph{weights at most \(w\)} if \(\mathcal H^i(K)\) has weights at most \(w+i\) for all \(i\), and it is \emph{pure of weight \(w\)} if each nonzero \(\mathcal H^i(K)\) is pure of weight \(w+i\).
\end{bpstatement}

% text: prose-13-04
\phantomsection\label{bp8}%
Deligne calls a complex \(K\) pure of weight \(w\) if \(K\) has weights at most \(w\) and its Verdier dual \(\mathbb D K\) has weights at most \(-w\). Over a point the two notions agree. In general neither implies the other, and we use ours only for sheaves and over a point.

The cohomology groups of Theorem \hyperref[bp36]{13.11} are Weil sheaves on \(\operatorname{Spec}k_0\). For them, having weights at most \(w\) is just a condition on the eigenvalues of Frobenius. Indeed, by Lemma \hyperref[bp9]{13.2} every Weil \(\overline{\mathbf Q}_\ell\)-sheaf over a point is mixed, and Lemma \hyperref[bp12]{13.3} reads the bound of a mixed sheaf off its eigenvalues.

% object: lemma-weil-sheaves-point-mixed (blueprint/chap-Weights/sec-Weil-sheaves-and-weights/lemma-weil-sheaves-point-mixed/_section.md)
\phantomsection\label{bp9}%
\begin{bpstatementit}{Lemma~13.2}{ (Weil sheaves on a point are mixed)}
Every constructible Weil \(\overline{\mathbf Q}_\ell\)-sheaf on \(\operatorname{Spec}k_0\) is mixed. A constructible Weil \(E\)-sheaf on \(\operatorname{Spec}k_0\) becomes mixed after extending scalars to a finite extension of \(E\) containing the eigenvalues of Frobenius.
\end{bpstatementit}

% proof: proof/lemma-weil-sheaves-point-mixed
\begin{proof}[{Proof}]\phantomsection\label{bp10}%
A constructible Weil sheaf on \(\operatorname{Spec}k_0\) is a finite-dimensional vector space \(V\) with an automorphism \(F\), the geometric Frobenius. In both cases we may assume that \(V\) is defined over a finite extension \(E'\) of \(\mathbf Q_\ell\) containing the eigenvalues of \(F\). For \(w\in\mathbf R\) let \(V_w\subset V\) be the sum of the generalized eigenspaces of \(F\) for the eigenvalues \(\alpha\) with \(|\iota(\alpha)|=q^{w/2}\). These subspaces are \(F\)-stable, only finitely many are nonzero, and \(V=\bigoplus_wV_w\). Let \(w_1<\dots<w_m\) be the weights that occur. Then the subspaces \(V_{w_1}\oplus\dots\oplus V_{w_j}\) form a finite filtration by Weil subsheaves with quotients \(V_{w_j}\), and \(V_{w_j}\) is pure of weight \(w_j\).
\end{proof}

% text: prose-13-05
\phantomsection\label{bp11}%
Extending scalars is needed for \(E\)-sheaves. Suppose \(2\) is not a square in \(\mathbf Q_\ell\). Then the rank-two \(\mathbf Q_\ell\)-sheaf on which Frobenius has eigenvalues \(1\pm\sqrt2\) is irreducible. Its two eigenvalues have different weights, so it has no filtration with pure quotients over \(\mathbf Q_\ell\).

% object: lemma-bounding-weights-mixed-sheaf (blueprint/chap-Weights/sec-Weil-sheaves-and-weights/lemma-bounding-weights-mixed-sheaf/_section.md)
\phantomsection\label{bp12}%
\begin{bpstatementit}{Lemma~13.3}{ (Bounding the weights of a mixed sheaf)}
A mixed sheaf has weights at most \(w\) if and only if it is punctually of weights \(\le w\).
\end{bpstatementit}

% proof: proof/lemma-bounding-weights-mixed-sheaf
\begin{proof}[{Proof}]\phantomsection\label{bp13}%
At every closed point \(x\) the eigenvalues of \(\operatorname{Frob}_x\) on the stalk of a mixed sheaf are those on the stalks of the quotients of its filtration. If the quotients have weights at most \(w\), the sheaf is therefore punctually of weights \(\le w\). Conversely, each nonzero quotient has a nonzero stalk at some closed point \(x\), and the eigenvalues of \(\operatorname{Frob}_x\) there are eigenvalues on the stalk of the sheaf. So if the sheaf is punctually of weights \(\le w\), the weight of the quotient is at most \(w\).
\end{proof}

% text: prose-13-06
\phantomsection\label{bp14}%
The lemma does not say that a sheaf punctually of weights \(\le w\) is mixed; the filtration has to be found. Nor need the filtration split: weights see only the characteristic polynomials of Frobenius, so a Jordan block and its semisimplification have the same weights. The following example is a mixed sheaf that is not pure and whose filtration does not split.

% object: example-kummer-extension (blueprint/chap-Weights/sec-Weil-sheaves-and-weights/example-kummer-extension/_section.md)
\phantomsection\label{bp15}%
\begin{bpstatement}{Example~13.4}{ (The Kummer extension)}
On \(\mathbf G_{m,k_0}\) with coordinate \(t\), the \(\mu_{\ell^n}\)-torsors \(u^{\ell^n}=t\) give the Kummer cocycle \(\kappa\), a continuous map from \(\pi_1(\mathbf G_{m,k_0})\) to \(\mathbf Z_\ell(1)\). Let \(\chi\) be the cyclotomic character, by which \(\pi_1\) acts on \(\mathbf Z_\ell(1)\). The matrices

\[
\sigma\longmapsto\begin{pmatrix}\chi(\sigma)&\kappa(\sigma)\\0&1\end{pmatrix}
\]

define a lisse sheaf \(\mathcal E\) on \(\mathbf G_{m,k_0}\), an extension \(0\to\overline{\mathbf Q}_\ell(1)\to\mathcal E\to\overline{\mathbf Q}_\ell\to0\). Geometric Frobenius acts on \(\overline{\mathbf Q}_\ell(1)\) by \(q_x^{-1}\). Hence at a closed point \(x\) the eigenvalues of \(\operatorname{Frob}_x\) are \(q_x^{-1}\) and \(1\), of weights \(-2\) and \(0\), and \(\mathcal E\) is mixed of weights at most \(0\) and not pure. The filtration does not split. Indeed, on \(\pi_1(\mathbf G_{m,k})\) the character \(\chi\) is trivial, and \(\kappa\) is a homomorphism onto \(\mathbf Z_\ell(1)\) because the covers \(u^{\ell^n}=t\) of \(\mathbf G_{m,k}\) are connected. So the line \(\overline{\mathbf Q}_\ell(1)\) is the only stable line. At each closed point, by contrast, the two eigenvalues are distinct and the stalk is the direct sum of its eigenlines.
\end{bpstatement}

% text: prose-13-07
\phantomsection\label{bp16}%
The behaviour of weights under the usual operations comes down to the eigenvalues of an operator on a filtered space, on duals and on tensor and exterior powers.

% object: lemma-filtrations-spectra (blueprint/chap-Weights/sec-Weil-sheaves-and-weights/lemma-filtrations-spectra/_section.md)
\phantomsection\label{bp17}%
\begin{bpstatementit}{Lemma~13.5}{ (Filtrations and spectra)}
Let \(V\) have a finite filtration preserved by an automorphism \(u\). Then

\[
\det(1-Tu\mid V)=\prod_a\det(1-Tu\mid\operatorname{gr}_aV).
\]

In particular, the eigenvalues, with algebraic multiplicity, are the union of the eigenvalues on the graded pieces. If the eigenvalues of \(u\) on \(V\) are \(\alpha_1,\ldots,\alpha_d\), then those on \(V^\vee\) are \(\alpha_i^{-1}\), and those on \(\bigwedge^rV\) are \(\alpha_{i_1}\cdots\alpha_{i_r}\) for \(i_1<\cdots<i_r\). If \(v\) on \(V'\) has eigenvalues \(\beta_j\), those of \(u\otimes v\) are \(\alpha_i\beta_j\).
\end{bpstatementit}

% proof: proof/lemma-filtrations-spectra
\begin{proof}[{Proof}]\phantomsection\label{bp18}%
In a basis adapted to the filtration, the matrix of \(u\) is block upper triangular, so its characteristic polynomial is the product of those of the diagonal blocks. For the other constructions we may choose bases over \(\overline{\mathbf Q}_\ell\) in which \(u\) and \(v\) are upper triangular. Then, in the induced tensor basis and in the ordered wedge basis, the resulting matrices are upper triangular with the stated diagonal entries. Finally, the matrix on the dual is the inverse transpose.
\end{proof}

% text: prose-13-08
\phantomsection\label{bp19}%
At a closed point, each operation in the next lemma and proposition acts on stalks by one of the constructions of the lemma: a subobject, a quotient, a tensor product or a dual, or, for an inverse image, a power of Frobenius. So the behaviour of weights is a matter of eigenvalues. Mixedness needs in addition some bookkeeping to produce the filtrations.

% object: lemma-subobjects-quotients-weil-sheaves (blueprint/chap-Weights/sec-Weil-sheaves-and-weights/lemma-subobjects-quotients-weil-sheaves/_section.md)
\phantomsection\label{bp20}%
\begin{bpstatementit}{Lemma~13.6}{ (Subobjects and quotients of Weil sheaves)}
Subobjects and quotients of a pure Weil sheaf of weight \(w\) are pure of weight \(w\). Subquotients of a mixed sheaf are mixed, and a subquotient of a sheaf that is mixed of weights at most \(w\) is mixed of weights at most \(w\).
\end{bpstatementit}

% proof: proof/lemma-subobjects-quotients-weil-sheaves
\begin{proof}[{Proof}]\phantomsection\label{bp21}%
By Lemma \hyperref[bp17]{13.5}, the eigenvalues on a subobject or a quotient occur among those of the whole. Hence subobjects and quotients of pure sheaves are pure. For mixed sheaves, we intersect a mixed filtration with a subobject, or push it into a quotient. The graded pieces are then subquotients of the original pure pieces, hence pure of the same weights, and so the upper bound persists.
\end{proof}

% object: proposition-weights-under-usual-operations (blueprint/chap-Weights/sec-Weil-sheaves-and-weights/proposition-weights-under-usual-operations/_section.md)
\phantomsection\label{bp22}%
\begin{bpstatementit}{Proposition~13.7}{ (Weights under the usual operations)}
Extensions of mixed sheaves are mixed, and an extension of two sheaves with a common upper bound has the same bound.  Tensor products of pure sheaves of weights \(a\) and \(b\) are pure of weight \(a+b\), and the dual of a pure lisse sheaf of weight \(a\) is pure of weight \(-a\). A tensor product of mixed sheaves of weights at most \(a\) and at most \(b\) is mixed of weights at most \(a+b\). The dual of a mixed lisse sheaf is mixed, of the negated weights; when its filtration has lisse graded pieces, their duals are the graded pieces of the annihilator filtration of the dual.  Inverse image preserves purity, mixedness and their upper bounds.  If a complex \(K\) has weights at most \(w\), then \(K[r]\) has weights at most \(w+r\).

For \(b\in\overline{\mathbf Q}_\ell^\times\), let \(\mathcal B_b\) be the geometrically constant rank-one Weil sheaf on which a degree-one Frobenius acts by \(b\). Tensoring by \(\mathcal B_b\) adds \(\lambda_\iota(b)=2\log_q|\iota(b)|\) to every weight.  The Tate sheaf \(\overline{\mathbf Q}_\ell(1)\) has geometric Frobenius \(q^{-1}\), so twisting by \((r)\) subtracts \(2r\) from every weight.
\end{bpstatementit}

% proof: proof/proposition-weights-under-usual-operations
\begin{proof}[{Proof}]\phantomsection\label{bp23}%
Each assertion has two parts: a statement about eigenvalues at closed points, which follows from Lemma \hyperref[bp17]{13.5}, and, for mixedness, the construction of a filtration with pure quotients.

An extension is filtered by a filtration of its subobject followed by the inverse images of a filtration of its quotient. The graded pieces of the result are those of the two given filtrations, so a common upper bound persists.

For tensor products we use that the coefficients form a field, so that the associated graded of the tensor product of two filtrations is the tensor product of the two associated gradeds. By Lemma \hyperref[bp17]{13.5}, the eigenvalues of tensor products and duals are the products \(\alpha_i\beta_j\) and the inverses \(\alpha_i^{-1}\), whose weights are \(a+b\) and \(-a\).

For the dual of a mixed lisse sheaf \(G\) with mixed filtration \((G_i)\), suppose first that the graded pieces \(G_i/G_{i-1}\) are lisse. Dualizing is exact on lisse sheaves, so the annihilator filtration \(G_i^\perp=\ker(G^\vee\to G_i^\vee)\) has graded pieces \((G_i/G_{i-1})^\vee\), which are pure of the negated weights by the computation above. In general the graded pieces need not be lisse: on a curve, a pure lisse \(G\) has the filtration \(j_!j^*G\subseteq G\), whose quotient \(i_*i^*G\) is a skyscraper with dual zero. We then argue by Noetherian induction on \(X\). Let \(U\subseteq X\) be a dense open, stable under Frobenius, on which all graded pieces are lisse, and let \(i:Z\to X\) be its closed complement. On \(U\) the first case applies to \(j^*G\). Since \(G\) is lisse, \(i^*(G^\vee)\simeq(i^*G)^\vee\), and \(i^*G\) is a mixed lisse sheaf on \(Z\), so \(i^*(G^\vee)\) is mixed by induction, of the negated weights. Extension by zero and direct image along a closed immersion are exact and do not change stalks, so they carry mixed filtrations to mixed filtrations of the same weights, and the exact sequence

\[
0\to j_!j^*(G^\vee)\to G^\vee\to i_*i^*(G^\vee)\to 0
\]

shows that \(G^\vee\) is mixed, of the negated weights.

Inverse images do not change weights, because the weight is normalized by the degree. Indeed, at a closed point of an inverse image the Frobenius is a positive power of the Frobenius at the image point, and the residue field degree is multiplied by the same integer. For shifts, \(\mathcal H^i(K[r])=\mathcal H^{i+r}(K)\) is mixed of weights at most \(w+i+r\), which is the bound \(w+r\) for \(K[r]\).

The local eigenvalue of \(\mathcal B_b\) at \(x\) is \(b^{\deg(x)}\), of weight \(2\log_q|\iota(b)|\), and tensoring by \(\mathcal B_b\) multiplies every eigenvalue by it. For the Tate sheaf, recall that arithmetic Frobenius acts on roots of unity by raising to the \(q\)th power. Geometric Frobenius is its inverse, so it acts on \(\overline{\mathbf Q}_\ell(1)\) by \(q^{-1}\), of weight \(-2\).
\end{proof}


% object: sec-Lisse-sheaves-as-representations (blueprint/chap-Weights/sec-Lisse-sheaves-as-representations/_section.md), content hidden
\subsection*{\texorpdfstring{13.2\quad Lisse sheaves as representations}{13.2 Lisse sheaves as representations}}
\phantomsection\label{bp24}%
\addcontentsline{toc}{subsection}{\texorpdfstring{13.2\quad Lisse sheaves as representations}{13.2 Lisse sheaves as representations}}
% text: prose-13-09
\phantomsection\label{bp25}%
Let \(X\) be a connected normal scheme of finite type over a field and \(L\) a lisse \(E_X\)-module. Lemma \bpout{3.24} gives lattices in \(L\) only on the strata of a finite stratification, and the nodal curve of Section \bpout{3.8} shows that without normality a global lattice need not exist. On a normal \(X\), however, we find a lattice on all of \(X\), and as a consequence lisse sheaves are the same as continuous representations of the fundamental group. Here normality is used to show that the étale \(X\)-scheme representing lattices in \(L\) has a component that is finite étale and surjective. Once we have such a component, we argue as in Lemma \bpout{3.24}.

% object: lemma-geometric-coverings-normal-scheme (blueprint/chap-Weights/sec-Lisse-sheaves-as-representations/lemma-geometric-coverings-normal-scheme/_section.md)
\phantomsection\label{bp26}%
\begin{bpstatementit}{Lemma~13.8}{ (Geometric coverings of a normal scheme)}
Let \(X\) be a nonempty connected normal scheme of finite type over a field. Let \(p:Y\to X\) be separated and étale. Suppose that for every valuation ring \(V\), with fraction field \(K\), every commutative diagram

\[
\begin{CD}
\operatorname{Spec}K @>>> Y\\
@VVV @VVpV\\
\operatorname{Spec}V @>>> X
\end{CD}
\]

has a lift \(\operatorname{Spec}V\to Y\). Then every connected component of \(Y\) is finite étale and surjective over \(X\).
\end{bpstatementit}

% proof: proof/lemma-geometric-coverings-normal-scheme
\begin{proof}[{Proof}]\phantomsection\label{bp27}%
We compare a component \(Y_1\) with the normalization \(N\) of \(X\) in the function field of \(Y_1\). By Zariski's main theorem, \(Y_1\) is an open subscheme of \(N\), and by the valuative hypothesis this open subscheme contains every point of \(N\).

Since \(Y\) is étale over \(X\), it is locally noetherian and normal [Stacks, \href{https://stacks.math.columbia.edu/tag/033C}{033C}]. Hence its connected components are integral, open and closed [Stacks, \href{https://stacks.math.columbia.edu/tag/033N}{033N}]. Fix a nonempty component \(Y_1\). Since a valuative lift whose generic point lies in \(Y_1\) stays in \(Y_1\), the hypothesis also holds for \(Y_1\to X\).

The image of \(Y_1\) is nonempty and open, hence contains the generic point \(\eta\) of \(X\). The generic fibre of \(Y_1\) is integral and étale over \(k(X)\), so it is \(\operatorname{Spec}k(Y_1)\). Thus the function field \(k(Y_1)\) is a finite separable extension of \(k(X)\), and we let \(N\) be the normalization of \(X\) in \(k(Y_1)\). Since \(X\) is Nagata, \(N\) is finite over \(X\) [Stacks, \href{https://stacks.math.columbia.edu/tag/035B}{035B}, \href{https://stacks.math.columbia.edu/tag/0AVK}{0AVK}]. We claim that \(Y_1\to X\) factors through a map \(f:Y_1\to N\). Indeed, since \(Y_1\) is normal, an element of \(k(Y_1)\) that is integral over the coordinate ring of an affine open of \(X\) lies in every local ring of \(Y_1\) over that open. These regular functions agree in \(k(Y_1)\), and hence they glue. The map \(f\) is separated, locally of finite type and locally quasi-finite, because its fibres lie in the zero-dimensional fibres of \(Y_1\to X\).

We show that \(f\) is an open immersion. Choose affine opens \(W\subseteq Y_1\) and \(N_0\subseteq N\) with \(f(W)\subseteq N_0\). By Zariski's main theorem [Stacks, \href{https://stacks.math.columbia.edu/tag/05K0}{05K0}], \(W\) is an open subscheme of a finite \(N_0\)-scheme, and the reduced closure of \(W\) in it is integral, finite and birational over \(N_0\), with function field \(k(Y_1)\). Since \(N_0\) is normal, this closure is \(N_0\) [Stacks, \href{https://stacks.math.columbia.edu/tag/0AB1}{0AB1}], and therefore \(W\to N_0\) is an open immersion. These local inverses glue: on the intersection of two such open images, the inverse sections agree generically, hence everywhere by separatedness, and so they glue to an inverse on the union of the images.

Finally, let \(n\in N\) and choose a valuation ring \(V\subseteq k(Y_1)\) dominating \(\mathcal O_{N,n}\) [Stacks, \href{https://stacks.math.columbia.edu/tag/00IA}{00IA}]. By hypothesis, the generic map \(\operatorname{Spec}k(Y_1)\to Y_1\) extends to a map \(\operatorname{Spec}V\to Y_1\), whose composite with \(f\) agrees with the canonical map \(\operatorname{Spec}V\to N\) generically and over \(X\). By separatedness of \(N\to X\) the two maps agree, so \(n\) lies in the open image of \(Y_1\). We conclude that \(Y_1=N\) is finite étale over \(X\). In particular its image is open and closed, and being nonempty, it is all of \(X\).
\end{proof}

% text: prose-13-10
\phantomsection\label{bp28}%
We apply the lemma to the scheme of lattices in a lisse sheaf. There the valuative hypothesis says that a lattice defined at the generic point of a valuation ring extends over the whole ring. This holds because, locally, lattices form a constant sheaf on a discrete set.

% object: lemma-global-lattice-normal-scheme (blueprint/chap-Weights/sec-Lisse-sheaves-as-representations/lemma-global-lattice-normal-scheme/_section.md)
\phantomsection\label{bp29}%
\begin{bpstatementit}{Lemma~13.9}{ (Global lattices on a normal scheme)}
Let \(X\) be a nonempty connected normal scheme of finite type over a field, and let \(E/\mathbf Q_\ell\) be finite. Every lisse \(E_X\)-module \(L\), in the pro-étale coefficient category of Section \bpout{2.12}, contains a lisse \(\mathcal O_{E,X}\)-submodule \(M\) with \(M[1/\ell]=L\). The reductions \(M_n=M/\pi_E^nM\) are finite locally constant étale sheaves, and \(M\simeq\varprojlim_nM_n\).
\end{bpstatementit}

% proof: proof/lemma-global-lattice-normal-scheme
\begin{proof}[{Proof}]\phantomsection\label{bp30}%
The scheme of lattices satisfies the valuative condition of Lemma \hyperref[bp26]{13.8}, so it has a finite étale surjective component, and we push the universal lattice of this component down to \(X\) by the trace.

We may assume that \(L\) has positive rank \(r\). Let \(Y\) be the \(X\)-scheme representing lattices in \(L\) (Lemma \bpout{3.15}), a separated surjective étale \(X\)-scheme with a universal lattice. Pro-étale locally on \(X\), the scheme \(Y\) is constant: the proof of that lemma shows that, for a quasi-compact pro-étale covering \(U\to X\), with \(U\) a pro-étale affine on which \(L\) is free,

\[
Y\times_XU\simeq
\coprod_{\GL_r(E)/\GL_r(\mathcal O_E)} U,
\]

where \(\GL_r(E)/\GL_r(\mathcal O_E)\) is the discrete set of lattices in \(E^r\).

We check the valuative property of Lemma \hyperref[bp26]{13.8}. Let \(V\) be a valuation ring with fraction field \(K\), and suppose given \(\operatorname{Spec}V\to X\) and a lift \(\xi:\operatorname{Spec}K\to Y\). Since \(U\to X\) is flat and surjective, there are a local extension of valuation rings \(V\to V'\) and a compatible map \(\operatorname{Spec}V'\to U\) [Stacks, \href{https://stacks.math.columbia.edu/tag/089F}{089F}]. This extension is faithfully flat, because it is local and \(V'\) is torsion-free, hence flat, over \(V\) [Stacks, \href{https://stacks.math.columbia.edu/tag/0539}{0539}]. By the displayed description of \(Y\times_XU\), the pullback of \(\xi\) to the spectrum of the fraction field of \(V'\) picks out a lattice in \(E^r\), and the constant section with this value extends it to a map \(\xi':\operatorname{Spec}V'\to Y\).

We descend \(\xi'\) by fpqc descent, which needs its two pullbacks to \(R=V'\otimes_VV'\) to agree. Over the domain \(V'\otimes_VK\), the maps \(\xi'\) and \(\xi\) agree at the generic point, hence everywhere by separatedness, so the two pullbacks agree after base change to \(K\). This suffices: since \(R\) is flat over \(V\), the map \(R\to R\otimes_VK\) is injective, and therefore \(\operatorname{Spec}(R\otimes_VK)\) is schematically dense in \(\operatorname{Spec}R\). As the equalizer of the two pullbacks is closed by separatedness, they agree on \(\operatorname{Spec}R\). Descent then gives a lift \(\operatorname{Spec}V\to Y\) restricting to \(\xi\), which is unique by separatedness and the density of \(\operatorname{Spec}K\).

By Lemma \hyperref[bp26]{13.8}, every nonempty component \(p_1:Y_1\to X\) of \(Y\) is therefore finite étale and surjective. We fix one, and let \(N\subseteq p_1^*L\) be its universal lattice. As in Lemma \bpout{3.24}, set

\[
M=\operatorname{im}\bigl(
p_{1*}N\longrightarrow p_{1*}p_1^*L
\xrightarrow{\operatorname{Tr}}L\bigr).
\]

The sheaf \(M\) is lisse with \(M[1/\ell]=L\). Indeed, after the cover is split into \(d\) sheets, this image is the sum of \(d\) lattices. On a common pro-étale refinement these become constant lattices in \(E^r\), and their sum is a finite free \(\mathcal O_E\)-lattice spanning \(E^r\).

By Theorem \bpout{2.43} the finite reductions \(M_n=M/\pi_E^nM\) are finite locally constant étale sheaves, and \(M=\varprojlim_nM_n\) by Lemma \bpout{2.51}, which we may check where \(M\) is free.
\end{proof}

% text: prose-13-11
\phantomsection\label{bp31}%
Galois theory of finite étale covers turns statements about the finite reductions of such an \(M\) into statements about representations.

% object: proposition-lattices-representations-normal-scheme (blueprint/chap-Weights/sec-Lisse-sheaves-as-representations/proposition-lattices-representations-normal-scheme/_section.md)
\phantomsection\label{bp32}%
\begin{bpstatementit}{Proposition~13.10}{ (Lattices and representations on a normal scheme)}
Let \(X\) be a nonempty connected normal scheme of finite type over a field, let \(E/\mathbf Q_\ell\) be finite, and fix a geometric point \(\bar x\). By Lemma \hyperref[bp29]{13.9}, every lisse \(E_X\)-module \(L\) contains a lisse \(\mathcal O_{E,X}\)-submodule \(M\) with \(M[1/\ell]=L\) and \(M\simeq\varprojlim_nM/\pi_E^nM\), whose reductions are finite locally constant étale sheaves. The fibre functor at \(\bar x\), \(L\mapsto\bigl(\varprojlim_n(M/\pi_E^nM)_{\bar x}\bigr)\otimes_{\mathcal O_E}E\), does not depend on the choice of \(M\) up to canonical isomorphism, and it gives an equivalence between lisse \(E_X\)-modules and finite-dimensional continuous \(E\)-representations of the ordinary profinite group \(\pi_1(X,\bar x)\).
\end{bpstatementit}

% proof: proof/proposition-lattices-representations-normal-scheme
\begin{proof}[{Proof}]\phantomsection\label{bp33}%
Let \(M\subset L\) be as in the statement, with reductions \(M_n=M/\pi_E^nM\). The fibres of these reductions define a continuous action of \(\pi_1(X,\bar x)\) on the finite free completed fibre

\[
M_{\bar x}:=\varprojlim_n(M_n)_{\bar x}.
\]

The fibre functor of the proposition takes \(L\) to the rational fibre \(M_{\bar x}\otimes_{\mathcal O_E}E\); in other words, we apply the fibre functor of finite covers to the finite reductions and complete the result.

Full faithfulness comes from the Galois theory of finite étale covers of the connected scheme \(X\). It says that taking the fibre at \(\bar x\) is an equivalence between finite locally constant sheaves on \(X_{\acute{e}t}\) and finite sets with continuous \(\pi_1(X,\bar x)\)-action, compatible with module structures. To pass from the finite reductions to \(L\) we need to clear denominators. Let \(M\subseteq L\) and \(M'\subseteq L'\) be lattices in two lisse \(E_X\)-modules. Then a morphism \(\phi:L\to L'\) has a uniform denominator: \(\pi_E^c\phi\) maps \(M\) into \(M'\) for some \(c\geq0\). To see this, bound the finitely many matrix entries on a quasi-compact pro-étale covering trivializing both lattices.

For faithfulness, suppose the fibre of \(\phi\) vanishes. Then the fibres of all the reductions of \(\pi_E^c\phi\) vanish, so these reductions vanish by the Galois correspondence. Since \(M'=\varprojlim_nM'_n\), it follows that \(\pi_E^c\phi=0\). For fullness, take an equivariant map between the rational fibres. After multiplication by some \(\pi_E^c\), it maps \(M_{\bar x}\) into \(M'_{\bar x}\). Its reductions are then equivariant maps of finite fibres, so by the Galois correspondence they come from compatible maps \(M_n\to M'_n\). The inverse limit of these is a map \(M\to M'\) whose rationalization, divided by \(\pi_E^c\), has the given fibre.

The same uniform denominator, applied to the identity of \(L\), shows that the rational fibre does not depend on the lattice. Indeed, if \(M_1\) and \(M_2\) are two lattices in \(L\), then \(\pi_E^c\,\mathrm{id}_{L}\) sends \(M_1\) into \(M_2\) for some \(c\geq0\), and symmetrically. On fibres these maps, divided by \(\pi_E^c\), therefore give mutually inverse equivariant isomorphisms between the two rational fibres.

For essential surjectivity we build a stable lattice and glue its reductions. A continuous representation of the profinite group \(\pi_1(X,\bar x)\) on \(E^r\) has compact image. Hence the \(\mathcal O_E\)-span of the translates of a lattice is bounded and contains that lattice. As \(\mathcal O_E\) is noetherian, this span is a stable finite free lattice, and its reductions define compatible finite locally constant étale sheaves. These sheaves become constant simultaneously on the surjective pro-étale cover obtained as the inverse limit of the finite étale covers corresponding to their kernels. Their inverse limit is therefore lisse over \(\mathcal O_E\), and its rationalization is lisse over \(E\). Finally, this construction is inverse to the fibre functor.
\end{proof}


% object: sec-The-Deligne-Laumon-theorem-2 (blueprint/chap-Weights/sec-The-Deligne-Laumon-theorem-2/_section.md), content hidden
\subsection*{\texorpdfstring{13.3\quad The Deligne–Laumon theorem}{13.3 The Deligne–Laumon theorem}}
\phantomsection\label{bp34}%
\addcontentsline{toc}{subsection}{\texorpdfstring{13.3\quad The Deligne–Laumon theorem}{13.3 The Deligne–Laumon theorem}}
% text: prose-13-12
\phantomsection\label{bp35}%
For a separated \(k_0\)-scheme of finite type \(X_0\), we write \(D_{\leq w}(X_0)\) for the category of bounded constructible Weil complexes \(K\) on \(X_0\) of weights at most \(w\) in the sense of Definition \hyperref[bp7]{13.1}, that is, such that \(\mathcal H^i(K)\) is mixed of weights at most \(w+i\) for every \(i\). The Deligne–Laumon theorem says that for a separated morphism of finite type \(f:X_0\to Y_0\), the functor \(Rf_!\) maps \(D_{\leq w}(X_0)\) into \(D_{\leq w}(Y_0)\). The shift by \(i\) makes this one statement for all degrees: for the structure morphism of \(X_0\), one condition on \(R\Gamma_c(X,K)\) bounds all the groups \(H^n_c(X,K)\) at once. Moreover, by Proposition \hyperref[bp22]{13.7}, replacing \(K\) by \(K[r]\) only moves the bound by \(r\). The proof occupies Chapter \bpout{18} and, for a single curve, the sections before it, and Corollary \bpout{19.2} deduces from the theorem the purity of the cohomology of a smooth proper variety.

% object: theorem-weights-rf-shriek (blueprint/chap-Weights/sec-The-Deligne-Laumon-theorem-2/theorem-weights-rf-shriek/_section.md)
\phantomsection\label{bp36}%
\begin{bpstatementit}{Theorem~13.11}{ (Weights of $Rf_!$)}
For every separated morphism of finite type \(f:X_0\to Y_0\),

\[
Rf_!D_{\leq w}(X_0)\subset D_{\leq w}(Y_0).
\]

In particular, if \(K\in D_{\leq w}(X_0)\), then \(H_c^n(X,K)\) has weights at most \(w+n\).
\end{bpstatementit}



\clearpage\phantomsection\hypertarget{bp-outside}{}%
\section*{Outside this export}
\addcontentsline{toc}{section}{Outside this export}
References marked \ensuremath{{}^\ast} are to parts of the blueprint this document leaves out:
\begin{description}
\item[Theorem 2.43] Comparison for constructible sheaves, in Ⅰ The six operations (\bpslug{theorem-{}\allowbreak{}comparison-{}\allowbreak{}constructible-{}\allowbreak{}sheaves})
\item[Section 2.12] Coefficient rings, in Ⅰ The six operations (\bpslug{sec-{}\allowbreak{}Coefficient-{}\allowbreak{}rings})
\item[Lemma 2.51] Comparison with completed coefficients, in Ⅰ The six operations (\bpslug{lemma-{}\allowbreak{}comparison-{}\allowbreak{}completed-{}\allowbreak{}coefficients})
\item[Lemma 3.15] Lattices in lisse $E$-sheaves, in Ⅰ The six operations (\bpslug{lemma-{}\allowbreak{}lattices-{}\allowbreak{}lisse-{}\allowbreak{}e-{}\allowbreak{}sheaves})
\item[Section 3.8] Lattices after stratification, in Ⅰ The six operations (\bpslug{sec-{}\allowbreak{}Lattices-{}\allowbreak{}after-{}\allowbreak{}stratification})
\item[Lemma 3.24] Lattices after a finite stratification, in Ⅰ The six operations (\bpslug{lemma-{}\allowbreak{}lattices-{}\allowbreak{}after-{}\allowbreak{}finite-{}\allowbreak{}stratification})
\item[Section 10.5] Weil complexes, in Ⅱ Duality and the trace formula (\bpslug{sec-{}\allowbreak{}Weil-{}\allowbreak{}complexes})
\item[Theorem 12.30] The trace formula as an Euler product, in Ⅱ Duality and the trace formula (\bpslug{theorem-{}\allowbreak{}trace-{}\allowbreak{}formula-{}\allowbreak{}euler-{}\allowbreak{}product})
\item[Theorem 14.5] Rank-one sheaves have finite geometric monodromy, in Ⅲ Weights and the Riemann hypothesis (\bpslug{theorem-{}\allowbreak{}rank-{}\allowbreak{}one-{}\allowbreak{}sheaves-{}\allowbreak{}have-{}\allowbreak{}finite})
\item[Section 14.2] Determinant weights, in Ⅲ Weights and the Riemann hypothesis (\bpslug{sec-{}\allowbreak{}Determinant-{}\allowbreak{}weights})
\item[Chapter 18] The Deligne–Laumon theorem in general, in Ⅲ Weights and the Riemann hypothesis (\bpslug{chap-{}\allowbreak{}The-{}\allowbreak{}Deligne-{}\allowbreak{}Laumon-{}\allowbreak{}theorem-{}\allowbreak{}in-{}\allowbreak{}general})
\item[Corollary 19.2] Purity of the cohomology of a smooth proper variety, in Ⅲ Weights and the Riemann hypothesis (\bpslug{corollary-{}\allowbreak{}purity-{}\allowbreak{}cohomology-{}\allowbreak{}smooth-{}\allowbreak{}proper-{}\allowbreak{}variety})
\item[Theorem 19.7] The Weil conjectures for smooth proper schemes, in Ⅲ Weights and the Riemann hypothesis (\bpslug{theorem-{}\allowbreak{}weil-{}\allowbreak{}conjectures-{}\allowbreak{}smooth-{}\allowbreak{}proper-{}\allowbreak{}schemes})
\end{description}
\clearpage\phantomsection%
\section*{References}
\addcontentsline{toc}{section}{References}
\bpbibitem{BBD82}{bp37}{A. Beilinson, J. Bernstein, and P. Deligne. \emph{Faisceaux pervers}, volume 100 of \emph{Astérisque}. 1982.}

\bpbibitem{BS15}{bp38}{B. Bhatt and P. Scholze. The pro-étale topology for schemes. \emph{Astérisque}, 369:99–201, 2015. \href{https://people.mpim-bonn.mpg.de/scholze/proetale.pdf}{Author's manuscript}.}

\bpbibitem{Con07}{bp39}{B. Conrad. Deligne's notes on Nagata compactifications. \emph{Journal of the Ramanujan Mathematical Society}, 22:205–257, 2007. \href{https://math.stanford.edu/\~conrad/papers/nagatafinal.pdf}{Author's manuscript}.}

\bpbibitem{Del74}{bp40}{P. Deligne. La conjecture de Weil. I. \emph{Publ. Math. IHÉS}, 43:273–307, 1974.}

\bpbibitem{Del80}{bp41}{P. Deligne. La conjecture de Weil. II. \emph{Publ. Math. IHÉS}, 52:137–252, 1980.}

\bpbibitem{Den22}{bp42}{C. Deninger. There is no “Weil-”cohomology theory with real coefficients for arithmetic curves. \emph{Ann. Sc. Norm. Super. Pisa Cl. Sci.}, 2022. To appear. arXiv:\href{https://arxiv.org/abs/2204.02714}{2204.02714}.}

\bpbibitem{Deu41}{bp43}{M. Deuring. Die Typen der Multiplikatorenringe elliptischer Funktionenkörper. \emph{Abh. Math. Sem. Univ. Hamburg}, 14:197–272, 1941.}

\bpbibitem{Eke90}{bp44}{T. Ekedahl. On the adic formalism. In \emph{The Grothendieck Festschrift, Vol. II}, volume 87 of \emph{Progress in Mathematics}, pages 197–218. Birkhäuser, 1990.}

\bpbibitem{Has36}{bp45}{H. Hasse. Zur Theorie der abstrakten elliptischen Funktionenkörper. I, II, III. \emph{J. Reine Angew. Math.}, 175:55–62, 69–88, 193–208, 1936.}

\bpbibitem{Hum72}{bp46}{J. E. Humphreys. \emph{Introduction to Lie Algebras and Representation Theory}, volume 9 of \emph{Graduate Texts in Mathematics}. Springer, 1972.}

\bpbibitem{Ill94}{bp47}{L. Illusie. Autour du théorème de monodromie locale. \emph{Astérisque}, 223:9–57, 1994. §4. \href{https://www.numdam.org/item/AST_1994__223__9_0/}{Original article}.}

\bpbibitem{Kat88}{bp48}{N. M. Katz. Travaux de Laumon. \emph{Astérisque}, 161–162:105–132, 1988. Séminaire Bourbaki, exp. 691.}

\bpbibitem{KW01}{bp49}{R. Kiehl and R. Weissauer. \emph{Weil conjectures, perverse sheaves and \(\ell\)-adic Fourier transform}, volume 42 of \emph{Ergebnisse der Mathematik}. Springer, 2001.}

\bpbibitem{Lau87}{bp50}{G. Laumon. Transformation de Fourier, constantes d'équations fonctionnelles et conjecture de Weil. \emph{Publ. Math. IHÉS}, 65:131–210, 1987.}

\bpbibitem{Mil86a}{bp51}{J. S. Milne. \href{https://www.jmilne.org/math/xnotes/AVs.pdf}{Abelian Varieties}. In \emph{Arithmetic Geometry}. 1986. Corrected online version.}

\bpbibitem{Mil86b}{bp52}{J. S. Milne. \href{https://www.jmilne.org/math/xnotes/JVs.pdf}{Jacobian Varieties}. In \emph{Arithmetic Geometry}. 1986. Corrected online version.}

\bpbibitem{Ran39}{bp53}{R. A. Rankin. Contributions to the theory of Ramanujan's function \(\tau(n)\) and similar arithmetical functions. II. \emph{Proc. Cambridge Philos. Soc.}, 35:357–372, 1939.}

\bpbibitem{SGA4}{bp54}{M. Artin, A. Grothendieck, and J.-L. Verdier, editors. \emph{Théorie des topos et cohomologie étale des schémas}, volumes 269, 270, 305 of \emph{Lecture Notes in Mathematics}. Springer, 1972–1973. SGA 4.}

\bpbibitem{SGA4½}{bp55}{P. Deligne. \emph{Cohomologie étale}, volume 569 of \emph{Lecture Notes in Mathematics}. Springer, 1977. SGA 4½; \emph{Théorèmes de finitude}, §3, and \emph{Rapport sur la formule des traces}, §3, and \emph{Sommes trigonométriques}, §1.}

\bpbibitem{SGA5}{bp56}{A. Grothendieck, editor. \emph{Cohomologie \(\ell\)-adique et fonctions \(L\)}, volume 589 of \emph{Lecture Notes in Mathematics}. Springer, 1977. SGA 5.}

\bpbibitem{SGA7II}{bp57}{P. Deligne and N. Katz, editors. \emph{Groupes de monodromie en géométrie algébrique}, volume 340 of \emph{Lecture Notes in Mathematics}. Springer, 1973. SGA 7 II; Exposé XIII, (2.1.7).}

\bpbibitem{Stacks}{bp58}{The Stacks Project Authors. \href{https://stacks.math.columbia.edu}{The Stacks Project}. Individual results are cited by their stable tags in the text.}

\bpbibitem{Var05}{bp59}{Y. Varshavsky. Lefschetz–Verdier trace formula and a generalization of a theorem of Fujiwara. 2005. Version 2. arXiv:\href{https://arxiv.org/abs/math/0505564}{math/0505564}. \href{https://arxiv.org/abs/math/0505564v2}{Author's preprint}.}

\bpbibitem{Wei48a}{bp60}{A. Weil. \emph{Sur les courbes algébriques et les variétés qui s'en déduisent}. Hermann, Paris, 1948.}

\bpbibitem{Wei48b}{bp61}{A. Weil. \emph{Variétés abéliennes et courbes algébriques}. Hermann, Paris, 1948.}

\bpbibitem{Wei49}{bp62}{A. Weil. Numbers of solutions of equations in finite fields. \emph{Bull. Amer. Math. Soc.}, 55:497–508, 1949.}

\end{document}
